\documentclass[../../main.tex]{subfiles}% 注意这里的文件路径不能用 ./main.tex ,否则用latexmk编译子文件会报错
\graphicspath{{\subfix{./image/}}} % 指定图片目录,后续可以直接使用图片文件名
% 注意这里的文件路径不能用 ../../image/ ,否则用latexmk编译子文件会报错

% 例如:
% \begin{figure}[H]
% \centering
% \includegraphics[scale=0.3]{图.png}
% \caption{}
% \label{figure:图}
% \end{figure}
% 注意:上述\label{}一定要放在\caption{}之后,否则引用图片序号会只会显示??.

\begin{document}

\section{$E_3$中的标架族}

\begin{definition}
取$E^3$中固定右手单位正交标架$\{O;\boldsymbol{i},\boldsymbol{j},\boldsymbol{k}\}$，则$E^3$中任意右手标架$\{p;\boldsymbol{e}_1,\boldsymbol{e}_2,\boldsymbol{e}_3\}$可表示为
\begin{align}
\begin{pmatrix}
\overrightarrow{Op}\\
\boldsymbol{e}_1\\
\boldsymbol{e}_2\\
\boldsymbol{e}_3
\end{pmatrix}
=
\begin{pmatrix}
a_1 & a_2 & a_3\\
a_{11} & a_{12} & a_{13}\\
a_{21} & a_{22} & a_{23}\\
a_{31} & a_{32} & a_{33}
\end{pmatrix}
\begin{pmatrix}
\boldsymbol{i}\\
\boldsymbol{j}\\
\boldsymbol{k}
\end{pmatrix},
\label{eq:frame_rep}
\end{align}
满足行列式条件
\begin{align}
\begin{vmatrix}
a_{11} & a_{12} & a_{13}\\
a_{21} & a_{22} & a_{23}\\
a_{31} & a_{32} & a_{33}
\end{vmatrix}
>0.
\label{eq:det_pos}
\end{align}
记全体右手标架构成的集合为\textbf{活动标架空间}$\mathfrak{F}$，其可视为$\mathbb{R}^{12}$中满足\eqref{eq:det_pos}的区域，坐标为$(a_1,a_2,a_3,a_{11},a_{12},\dots,a_{33})$。

若标架$\{p;\boldsymbol{e}_1,\boldsymbol{e}_2,\boldsymbol{e}_3\}$为\textbf{右手单位正交标架}，则还满足正交条件
\begin{align}
\boldsymbol{e}_i\cdot\boldsymbol{e}_j=\delta_{ij},\quad 1\leqslant i,j\leqslant3,
\label{eq:ortho_cond}
\end{align}
即
\begin{align}
\begin{pmatrix}
a_{11} & a_{12} & a_{13}\\
a_{21} & a_{22} & a_{23}\\
a_{31} & a_{32} & a_{33}
\end{pmatrix}
\begin{pmatrix}
a_{11} & a_{21} & a_{31}\\
a_{12} & a_{22} & a_{32}\\
a_{13} & a_{23} & a_{33}
\end{pmatrix}
=
\begin{pmatrix}
1 & 0 & 0\\
0 & 1 & 0\\
0 & 0 & 1
\end{pmatrix}.
\label{eq:ortho_mat}
\end{align}
记全体右手单位正交标架构成的集合为$\tilde{\mathfrak{F}}$，其为$\mathbb{R}^{12}$中满足\eqref{eq:det_pos}与\eqref{eq:ortho_mat}的6维代数曲面。
\end{definition}

\begin{definition}
对标架$\{p;\boldsymbol{e}_1,\boldsymbol{e}_2,\boldsymbol{e}_3\}$，对\eqref{eq:frame_rep}微分得
\begin{align}
\begin{pmatrix}
\mathrm{d}\overrightarrow{Op}\\
\mathrm{d}\boldsymbol{e}_1\\
\mathrm{d}\boldsymbol{e}_2\\
\mathrm{d}\boldsymbol{e}_3
\end{pmatrix}
=
\begin{pmatrix}
\mathrm{d}a_1 & \mathrm{d}a_2 & \mathrm{d}a_3\\
\mathrm{d}a_{11} & \mathrm{d}a_{12} & \mathrm{d}a_{13}\\
\mathrm{d}a_{21} & \mathrm{d}a_{22} & \mathrm{d}a_{23}\\
\mathrm{d}a_{31} & \mathrm{d}a_{32} & \mathrm{d}a_{33}
\end{pmatrix}
\begin{pmatrix}
\boldsymbol{i}\\
\boldsymbol{j}\\
\boldsymbol{k}
\end{pmatrix}.
\label{eq:diff_frame_rep}
\end{align}
设$(b_{ij})=(a_{ij})^{-1}$，即
\begin{align}
\begin{pmatrix}
\boldsymbol{i}\\
\boldsymbol{j}\\
\boldsymbol{k}
\end{pmatrix}
=
\begin{pmatrix}
b_{11} & b_{12} & b_{13}\\
b_{21} & b_{22} & b_{23}\\
b_{31} & b_{32} & b_{33}
\end{pmatrix}
\begin{pmatrix}
\boldsymbol{e}_1\\
\boldsymbol{e}_2\\
\boldsymbol{e}_3
\end{pmatrix},
\label{eq:dual_basis}
\end{align}
代入\eqref{eq:diff_frame_rep}，展开得到\textbf{活动标架运动公式}：
\begin{align}
\mathrm{d}\overrightarrow{Op}&=\Omega^1\boldsymbol{e}_1+\Omega^2\boldsymbol{e}_2+\Omega^3\boldsymbol{e}_3,\nonumber\\
\mathrm{d}\boldsymbol{e}_i&=\Omega_i^1\boldsymbol{e}_1+\Omega_i^2\boldsymbol{e}_2+\Omega_i^3\boldsymbol{e}_3,\quad 1\leqslant i\leqslant3,
\label{eq:motion_formula}
\end{align}
其中
\begin{align}
\Omega^j=\sum\limits_{k=1}^3\mathrm{d}a_k\cdot b_{kj},\quad \Omega_i^j=\sum\limits_{k=1}^3\mathrm{d}a_{ik}\cdot b_{kj},\quad 1\leqslant i,j\leqslant3
\label{eq:rel_comp_def}
\end{align}
称为\textbf{活动标架的相对分量}，为$\mathfrak{F}$上的1次外微分式。

对单位正交标架，$(b_{ij})=(a_{ji})$，故相对分量为
\begin{align}
\Omega^j=\sum\limits_{k=1}^3a_{jk}\mathrm{d}a_k,\quad \Omega_i^j=\sum\limits_{k=1}^3a_{jk}\mathrm{d}a_{ik},\quad 1\leqslant i,j\leqslant3.
\label{eq:ortho_rel_comp}
\end{align}
\end{definition}

\begin{proposition}
单位正交活动标架的相对分量满足反对称性：
\begin{align}
\Omega_i^j=-\Omega_j^i,\quad 1\leqslant i,j\leqslant3.
\label{eq:antisym_omega}
\end{align}
因此独立的相对分量仅有6个：$\Omega^1,\Omega^2,\Omega^3,\Omega_1^2=-\Omega_2^1,\Omega_1^3=-\Omega_3^1,\Omega_2^3=-\Omega_3^2$。
\end{proposition}
\begin{proof}

由单位正交条件\eqref{eq:ortho_mat}的分量形式，对其微分即可验证$\Omega_i^j=-\Omega_j^i$。

\end{proof}

\begin{theorem}
活动标架的相对分量满足\textbf{标架空间结构方程}：
\begin{align}
\mathrm{d}\Omega^j=\sum\limits_{k=1}^3\Omega^k\wedge\Omega_k^j,\quad \mathrm{d}\Omega_i^j=\sum\limits_{k=1}^3\Omega_i^k\wedge\Omega_k^j,\quad 1\leqslant i,j\leqslant3.
\label{eq:struct_eq_global}
\end{align}
\end{theorem}
\begin{proof}

由外微分幂零性$\mathrm{d}\circ\mathrm{d}=0$，对运动公式\eqref{eq:motion_formula}两边取外微分：
\begin{align*}
0=\mathrm{d}(\mathrm{d}\overrightarrow{Op})&=\sum\limits_{j=1}^3\mathrm{d}(\Omega^j\boldsymbol{e}_j)=\sum\limits_{j=1}^3\left(\mathrm{d}\Omega^j-\sum\limits_{k=1}^3\Omega^k\wedge\Omega_k^j\right)\boldsymbol{e}_j,\\
0=\mathrm{d}(\mathrm{d}\boldsymbol{e}_i)&=\sum\limits_{j=1}^3\mathrm{d}(\Omega_i^j\boldsymbol{e}_j)=\sum\limits_{j=1}^3\left(\mathrm{d}\Omega_i^j-\sum\limits_{k=1}^3\Omega_i^k\wedge\Omega_k^j\right)\boldsymbol{e}_j.
\end{align*}
因$\boldsymbol{e}_1,\boldsymbol{e}_2,\boldsymbol{e}_3$线性无关，故系数必为零，即得结构方程。

\end{proof}

\begin{remark}
单位正交标架空间$\tilde{\mathfrak{F}}$的结构方程由反对称性\eqref{eq:antisym_omega}简化为：
\begin{align*}
\mathrm{d}\Omega^j&=\sum\limits_{k=1}^3\Omega^k\wedge\Omega_k^j,\\
\mathrm{d}\Omega_1^2&=\Omega_1^3\wedge\Omega_3^2=-\Omega_1^3\wedge\Omega_2^3,\\
\mathrm{d}\Omega_1^3&=\Omega_1^2\wedge\Omega_2^3,\\
\mathrm{d}\Omega_2^3&=\Omega_2^1\wedge\Omega_1^3=\Omega_1^2\wedge\Omega_3^2.
\end{align*}
\end{remark}

\begin{definition}
设$\tilde{D}\subset\mathbb{R}^r$为区域，\textbf{$r$参数标架族}为$C^\infty$映射$\sigma:\tilde{D}\to\mathfrak{F}$，即
\begin{align}
a_i=a_i(u^1,\dots,u^r),\quad a_{ij}=a_{ij}(u^1,\dots,u^r),\quad \det(a_{ij})>0.
\label{eq:frame_family}
\end{align}
对\eqref{eq:frame_family}微分，得到参数域$\tilde{D}$上的\textbf{诱导运动公式}：
\begin{align}
\mathrm{d}\overrightarrow{Op}=\sum\limits_{k=1}^3\omega^k\boldsymbol{e}_k,\quad \mathrm{d}\boldsymbol{e}_i=\sum\limits_{k=1}^3\omega_i^k\boldsymbol{e}_k,
\label{eq:induced_motion}
\end{align}
其中
\begin{align}
\omega^k=\sigma^*\Omega^k,\quad \omega_i^k=\sigma^*\Omega_i^k
\label{eq:induced_rel_comp}
\end{align}
为$\tilde{D}$上的1次外微分式，称为\textbf{标架族的相对分量}。

若$\sigma:\tilde{D}\to\tilde{\mathfrak{F}}$，则为\textbf{$r$参数单位正交标架族}，满足
\begin{align}
\sum\limits_{k=1}^3a_{ik}a_{jk}=\delta_{ij},
\label{eq:ortho_family_cond}
\end{align}
且诱导相对分量满足$\omega_i^j=-\omega_j^i$。
\end{definition}

\begin{theorem}
$r$参数标架族的诱导相对分量满足\textbf{诱导结构方程}：
\begin{align}
\mathrm{d}\omega^j=\sum\limits_{k=1}^3\omega^k\wedge\omega_k^j,\quad \mathrm{d}\omega_i^j=\sum\limits_{k=1}^3\omega_i^k\wedge\omega_k^j,\quad 1\leqslant i,j\leqslant3.
\label{eq:struct_eq_induced}
\end{align}
\end{theorem}
\begin{proof}

由外微分的拉回性质$\mathrm{d}\circ\sigma^*=\sigma^*\circ\mathrm{d}$，结合全局结构方程\eqref{eq:struct_eq_global}：
\begin{align*}
\mathrm{d}\omega^j&=\mathrm{d}(\sigma^*\Omega^j)=\sigma^*(\mathrm{d}\Omega^j)=\sigma^*\left(\sum\limits_{k=1}^3\Omega^k\wedge\Omega_k^j\right)=\sum\limits_{k=1}^3\omega^k\wedge\omega_k^j,\\
\mathrm{d}\omega_i^j&=\mathrm{d}(\sigma^*\Omega_i^j)=\sigma^*(\mathrm{d}\Omega_i^j)=\sigma^*\left(\sum\limits_{k=1}^3\Omega_i^k\wedge\Omega_k^j\right)=\sum\limits_{k=1}^3\omega_i^k\wedge\omega_k^j.
\end{align*}

\end{proof}

\begin{theorem}
任给$\tilde{D}\subset\mathbb{R}^r$上的12个1次外微分式$\omega^j,\omega_i^j\ (1\leqslant i,j\leqslant3)$，若满足诱导结构方程\eqref{eq:struct_eq_induced}，则存在$E^3$中依赖$r$个参数的右手标架族，以$\omega^j,\omega_i^j$为相对分量。
\end{theorem}

\begin{theorem}
任给$\tilde{D}\subset\mathbb{R}^r$上的6个反对称1次外微分式
\begin{align*}
\omega^1,\omega^2,\omega^3,\quad \omega_1^2=-\omega_2^1,\ \omega_1^3=-\omega_3^1,\ \omega_2^3=-\omega_3^2,
\end{align*}
若满足诱导结构方程\eqref{eq:struct_eq_induced}，则存在$E^3$中依赖$r$个参数的右手单位正交标架族，以$\omega^j,\omega_i^j$为相对分量；且任意两个该类标架族可通过$E^3$的刚体运动重合。
\end{theorem}

\begin{remark}
上述存在性定理等价于一阶线性偏微分方程组的完全可积性条件，证明可参见附录相关定理。
\end{remark}

\begin{example}
设$p\in E^3$为定点，以$p$为原点的全体右手标架构成依赖9个参数的标架族，拓扑等价于$\mathrm{GL}^+(3)$；求其相对分量。
\end{example}
\begin{solution}

定点条件为$a_i=\text{常数}\ (1\leqslant i\leqslant3)$，故
\begin{align*}
\omega^1=\omega^2=\omega^3=0,\quad \omega_i^j=\sum\limits_{k=1}^3b_{kj}\mathrm{d}a_{ik},
\end{align*}
其中$(b_{kj})=(a_{ij})^{-1}$。

对单位正交情形，拓扑等价于$\mathrm{SO}(3)$，相对分量为
\begin{align*}
\omega^1=\omega^2=\omega^3=0,\quad \omega_i^j=\sum\limits_{k=1}^3a_{jk}\mathrm{d}a_{ik},
\end{align*}
满足$\sum\limits_{k=1}^3a_{ik}a_{jk}=\delta_{ij},\ \det(a_{ij})=1$。

\end{solution}

\begin{example}
设$C\subset E^3$为挠率非零的正则曲线，弧长参数$s$，曲率$\kappa(s)$，挠率$\tau(s)$，其Frenet标架场$\{\boldsymbol{r}(s);\boldsymbol{\alpha}(s),\boldsymbol{\beta}(s),\boldsymbol{\gamma}(s)\}$为单参数单位正交标架族；求其相对分量。
\end{example}
\begin{solution}

由Frenet公式：
\begin{align*}
\mathrm{d}\boldsymbol{r}&=\boldsymbol{\alpha}\mathrm{d}s,\\
\mathrm{d}\boldsymbol{\alpha}&=\kappa\boldsymbol{\beta}\mathrm{d}s,\\
\mathrm{d}\boldsymbol{\beta}&=-\kappa\boldsymbol{\alpha}\mathrm{d}s+\tau\boldsymbol{\gamma}\mathrm{d}s,\\
\mathrm{d}\boldsymbol{\gamma}&=-\tau\boldsymbol{\beta}\mathrm{d}s,
\end{align*}
对比运动公式，得相对分量：
\begin{align*}
\omega^1=\mathrm{d}s,\ \omega^2=\omega^3=0,\quad \omega_1^1=\omega_2^2=\omega_3^3=0,\ \omega_1^2=-\omega_2^1=\kappa\mathrm{d}s,\quad \omega_1^3=-\omega_3^1=0,\ \omega_2^3=-\omega_3^2=\tau\mathrm{d}s.
\end{align*}

反之，给定$\kappa(s)>0,\tau(s)$，可构造对应1次外微分式，满足结构方程，故存在唯一（差刚体运动）的正则曲线，此即\textbf{曲线论基本定理}。

\end{solution}

\begin{example}
沿挠率非零曲线$C$，取$e_1$为切向量，$\{e_2,e_3\}$可绕$e_1$旋转，构成依赖2个参数$(s,\theta)$的单位正交标架族；求其相对分量，并给出曲率的标架分量表示。
\end{example}
\begin{solution}

设
\begin{align*}
\boldsymbol{e}_2&=\cos\theta\boldsymbol{\beta}+\sin\theta\boldsymbol{\gamma},\\
\boldsymbol{e}_3&=-\sin\theta\boldsymbol{\beta}+\cos\theta\boldsymbol{\gamma},
\end{align*}
计算微分得相对分量：
\begin{align*}
\omega^1=\mathrm{d}s,\ \omega^2=\omega^3=0,\quad \omega_1^1=\omega_2^2=\omega_3^3=0,\ \omega_1^2=-\omega_2^1=\kappa\cos\theta\mathrm{d}s,\quad \omega_1^3=-\omega_3^1=-\kappa\sin\theta\mathrm{d}s,\ \omega_2^3=-\omega_3^2=\mathrm{d}\theta+\tau\mathrm{d}s.
\end{align*}

取单参数一阶标架族（$\theta=\theta(s)$），则
\begin{align*}
(\omega_1^2)^2+(\omega_1^3)^2=(\kappa\mathrm{d}s)^2,
\end{align*}
故曲率的标架不变量表示为
\begin{align}
\kappa=\frac{\sqrt{(\omega_1^2)^2+(\omega_1^3)^2}}{\mathrm{d}s}.
\label{eq:curvature_frame}
\end{align}
Frenet标架为二阶标架族，对应$\theta=0$。

\end{solution}

\begin{example}
求$E^3$中球坐标系给出的自然标架场、对应的单位正交标架场，并求解其相对分量。
\end{example}
\begin{solution}

设$(x,y,z)$为$E^3$中的笛卡儿直角坐标系，球坐标变换为
\begin{align*}
x = r\cos\theta\cos\varphi,\quad y = r\cos\theta\sin\varphi,\quad z = r\sin\theta.
\end{align*}
球坐标系对应的\textbf{自然标架场}为$\left\{\bm{r};\frac{\partial \bm{r}}{\partial r},\frac{\partial \bm{r}}{\partial \varphi},\frac{\partial \bm{r}}{\partial \theta}\right\}$，其中
\begin{align}
\frac{\partial \bm{r}}{\partial r} &= (\cos\theta\cos\varphi,\cos\theta\sin\varphi,\sin\theta), \nonumber \\
\frac{\partial \bm{r}}{\partial \varphi} &= (-r\cos\theta\sin\varphi,r\cos\theta\cos\varphi,0), \nonumber \\
\frac{\partial \bm{r}}{\partial \theta} &= (-r\sin\theta\cos\varphi,-r\sin\theta\sin\varphi,r\cos\theta). \label{eq:natural-frame-sphere}
\end{align}
可验证该组自然标架两两正交且构成右手系。将其单位化，得到\textbf{单位正交标架场}$\{\bm{r}; \boldsymbol{e}_1,\boldsymbol{e}_2,\boldsymbol{e}_3\}$，其中
\begin{align}
\boldsymbol{e}_1 &= \frac{\partial \bm{r}}{\partial r} = (\cos\theta\cos\varphi,\cos\theta\sin\varphi,\sin\theta), \nonumber \\
\boldsymbol{e}_2 &= \frac{\partial \bm{r}}{\partial \varphi}\bigg/\left|\frac{\partial \bm{r}}{\partial \varphi}\right| = (-\sin\varphi,\cos\varphi,0), \nonumber \\
\boldsymbol{e}_3 &= \frac{\partial \bm{r}}{\partial \theta}\bigg/\left|\frac{\partial \bm{r}}{\partial \theta}\right| = (-\sin\theta\cos\varphi,-\sin\theta\sin\varphi,\cos\theta). \label{eq:ortho-frame-sphere}
\end{align}

对向径$\bm{r}$求微分：
\begin{align}
\mathrm{d}\bm{r} = \frac{\partial \bm{r}}{\partial r}\mathrm{d}r + \frac{\partial \bm{r}}{\partial \varphi}\mathrm{d}\varphi + \frac{\partial \bm{r}}{\partial \theta}\mathrm{d}\theta = \mathrm{d}r\,\boldsymbol{e}_1 + r\cos\theta\mathrm{d}\varphi\,\boldsymbol{e}_2 + r\mathrm{d}\theta\,\boldsymbol{e}_3. \label{eq:dr-sphere}
\end{align}
由此求得单位正交标架的对偶1-形式：
\begin{align}
\omega^1 &= \mathrm{d}\bm{r}\cdot\boldsymbol{e}_1 = \mathrm{d}r, \nonumber \\
\omega^2 &= \mathrm{d}\bm{r}\cdot\boldsymbol{e}_2 = r\cos\theta\mathrm{d}\varphi, \nonumber \\
\omega^3 &= \mathrm{d}\bm{r}\cdot\boldsymbol{e}_3 = r\mathrm{d}\theta. \label{eq:dual-form-sphere}
\end{align}

进一步计算标架的相对分量：
\begin{align}
\omega_1^2 &= -\omega_2^1 = \cos\theta\mathrm{d}\varphi, \nonumber \\
\omega_1^3 &= -\omega_3^1 = \mathrm{d}\theta, \nonumber \\
\omega_2^3 &= -\omega_3^2 = \sin\theta\mathrm{d}\varphi. \label{eq:rel-comp-sphere}
\end{align}

\end{solution}

\begin{example}
计算$(3.32)$式.
\end{example}

\begin{example}
求通过柱坐标系在$E^3$中建立的单位正交标架场，并且求它的相对分量.
\end{example}

\begin{example}
设$C$是$Oxy$平面上以$O$为圆心、以$R$为半径的圆周，在$\mathbb{R}^3$中以$C$为底的直圆柱面以外的点集$U=\{(x,y,z)\in\mathbb{R}^3:x^2+y^2>R^2\}$上引进如下的参数系$(\rho,\theta,\psi)$，使得
\begin{gather*}
x = (R+\rho\cos\psi)\cos\theta, \\
y = (R+\rho\cos\psi)\sin\theta, \\
z = \rho\sin\psi,\quad \rho>0,\ -\frac{\pi}{2}<\psi<\frac{\pi}{2},\ -\pi<\theta<\pi.
\end{gather*}
证明：参数系$(\rho,\theta,\psi)$给出的参数曲线网是正交曲线网. 在$U$上建立单位正交标架场$\{\bm{r};\bm{e}_1,\bm{e}_2,\bm{e}_3\}$，使得$\bm{e}_1,\bm{e}_2,\bm{e}_3$分别是$\rho$-曲线、$\theta$-曲线、$\psi$-曲线的切向量，求这个标架场的相对分量.
\end{example}

\begin{example}
设$f$是定义在$E^3$上的连续可微函数，命
\begin{align*}
\bm{e}_1 = \frac{\sqrt{2}}{2}(\sin f,1,-\cos f), \quad
\bm{e}_2 = \frac{\sqrt{2}}{2}(\sin f,-1,-\cos f), \quad
\bm{e}_3 = (-\cos f,0,-\sin f).
\end{align*}
验证$\{\bm{r};\bm{e}_1,\bm{e}_2,\bm{e}_3\}$是$E^3$中的单位正交标架场，并且求它的相对分量.
\end{example}




\end{document}